\documentclass[11pt,letterpaper]{article}

\usepackage{amsmath,amssymb,amsthm}
\usepackage{mathtools}
\usepackage{booktabs}
\usepackage{array}
\usepackage{geometry}
\usepackage{hyperref}
\usepackage{cleveref}
\usepackage{enumitem}
\usepackage{xcolor}
\usepackage{algorithm}
\usepackage{algpseudocode}

\geometry{margin=1.1in}

\hypersetup{colorlinks=true, linkcolor=blue!60!black, citecolor=blue!60!black}

\newcommand{\dd}{\mathrm{d}}
\newcommand{\dt}{\Delta t}
\newcommand{\mol}{\,\text{mol}}
\newcommand{\J}{\,\text{J}}
\newcommand{\W}{\,\text{W}}
\newcommand{\K}{\,\text{K}}
\newcommand{\Pa}{\,\text{Pa}}
\newcommand{\mL}{\,\text{mL}}

\title{\textbf{D86 Distillation Simulation Algorithm}\\[0.5em]
\large \texttt{run\_d86\_simulation\_rk2\_condenser\_optE\_v2}\\[0.3em]
\normalsize Physics-corrected condenser model with Option~E controller}
\author{FuelLib}
\date{\today}

\begin{document}
\maketitle
\tableofcontents
\newpage

% ---------------------------------------------------------------------------
\section{Overview}
\label{sec:overview}

This document describes the equations and numerical algorithm implemented in
\texttt{run\_d86\_simulation\_rk2\_condenser\_optE\_v2}, the physics-corrected
D86 distillation simulation in FuelLib.  The model tracks the evolution of a
100\,mL fuel charge in an ASTM~D86 distillation apparatus, predicting the
temperature registered by the column thermocouple (the ``D86 temperature'')
as a function of distillate volume collected.

The simulation couples four physical stages in sequence at each time step:

\begin{enumerate}[nosep]
  \item \textbf{Stage~1} — Pot bubble-point temperature and vapour generation rate.
  \item \textbf{Stage~2} — Column neck partial condensation with natural-convection heat loss.
  \item \textbf{Stage~2b} — Lumped condenser tube with isothermal cooling bath.
  \item \textbf{Stage~3} — Thermocouple bulb thermal lag (first-order lumped capacitance).
\end{enumerate}

Time integration uses the second-order Runge--Kutta method (Heun's predictor--corrector).
An Option~E inverse-model feed-forward controller maintains the target distillation
rate of 4.5\,mL/min.

Five physics corrections relative to the baseline condenser model are incorporated;
they are flagged throughout as \textbf{[Issue~$n$]}.

% ---------------------------------------------------------------------------
\section{Notation and State Variables}
\label{sec:notation}

\subsection{Mixture description}

The fuel is a mixture of $n_c$ compounds.  The pot is described by the component
mole vector
\begin{equation}
  \mathbf{N} = (N_1, \ldots, N_{n_c})^\top \quad [\mol],
\end{equation}
with total moles $W = \sum_i N_i$ and pot mole fractions $X_i = N_i / W$.
Mass fractions are obtained via $Y_i = X_i M_i / \bar{M}$, where $M_i$
is the molar mass of compound~$i$ and $\bar{M} = \sum_i X_i M_i$.

\subsection{Key state variables carried between steps}

\begin{center}
\begin{tabular}{lll}
\toprule
Symbol & Units & Description \\
\midrule
$\mathbf{N}(t)$ & mol & Component moles in the pot \\
$T_2(t)$ & K & Column neck temperature (carried for energy balance) \\
$R_2(t)$ & mol/s & Reflux return rate from neck \\
$T_{\mathrm{D86}}(t)$ & K & Thermocouple (D86) temperature \\
$Q_1(t)$ & W & Heater power (updated by controller each step) \\
$\mathbf{c}_{\mathrm{cond}}(t)$ & --- & Condenser-outlet mole fractions (Issue~11) \\
\bottomrule
\end{tabular}
\end{center}

\subsection{Fixed parameters}

\begin{center}
\begin{tabular}{lll}
\toprule
Symbol & Units & Description \\
\midrule
$P$ & Pa & System pressure (typically 101\,325\,Pa) \\
$T_{\mathrm{room}}$ & K & Ambient temperature \\
$T_{\mathrm{bath}}$ & K & Condenser cooling-bath temperature \\
$UA_{\mathrm{cond}}$ & W/K & Condenser overall heat-transfer coefficient $\times$ area \\
$h_{\mathrm{mult}}$ & --- & Natural-convection scaling factor for neck \\
$C_{\mathrm{glass}}$ & J/K & Thermocouple bulb heat capacity \\
$\dot{V}_{\mathrm{target}}$ & mL/min & Target distillation rate (4.5\,mL/min) \\
$K_{p,\mathrm{trim}}$ & W\,min/mL & Proportional trim gain \\
$\dt$ & s & Integration time step \\
\bottomrule
\end{tabular}
\end{center}

% ---------------------------------------------------------------------------
\section{Vapour--Liquid Equilibrium}
\label{sec:vle}

\subsection{K-values (modified Raoult's law)}
\label{sec:kvalues}

Vapour--liquid K-values are computed from modified Raoult's law with UNIFAC
activity coefficients:
\begin{equation}
  K_i(T, P, \mathbf{x}) = \frac{\gamma_i(\mathbf{x}, T)\, P_i^{\mathrm{sat}}(T)}{P},
  \label{eq:kvalue}
\end{equation}
where $\gamma_i$ is the liquid-phase activity coefficient from UNIFAC and
$P_i^{\mathrm{sat}}$ is the pure-component vapour pressure from the Antoine
equation (both evaluated by \texttt{FuelLib.fuel}).

\subsection{Rachford--Rice isothermal flash}
\label{sec:rr}

Given a feed of composition $z_i$ at temperature $T$ and pressure $P$,
the two-phase Rachford--Rice equation is solved for the vapour fraction $V$:
\begin{equation}
  \sum_{i=1}^{n_c} \frac{z_i(K_i - 1)}{1 + V(K_i - 1)} = 0.
  \label{eq:rr}
\end{equation}
The solution interval is the Whitson--Brul\'e negative-flash bracket
$\bigl[1/(1-K_{\max}),\, 1/(1-K_{\min})\bigr]$, solved by \texttt{scipy.optimize.bisect}.

The outer loop is a successive-substitution (SS) update on the K-values:
starting from $K_i^{(0)}$ computed with the feed composition $z_i$, each
iteration solves~\cref{eq:rr} for $V$, then computes updated liquid compositions
\begin{equation}
  x_i = \frac{z_i}{1 + V(K_i - 1)},
  \qquad
  y_i = K_i x_i \big/ \textstyle\sum_j K_j x_j,
\end{equation}
and re-evaluates $K_i^{(k+1)}$ from \cref{eq:kvalue} using the updated $x_i$
(and $y_i$ as the vapour-phase guess for SRK).  Convergence is declared when
$\|K^{(k+1)} - K^{(k)}\|_\infty / \|K^{(k)}\|_\infty < 10^{-5}$.

Single-phase feeds are detected without entering the SS loop:
\begin{equation}
  \text{all liquid if } \sum_i z_i(K_i-1) \le 0; \qquad
  \text{all vapour if } \sum_i z_i\tfrac{K_i-1}{K_i} \ge 0.
\end{equation}

% ---------------------------------------------------------------------------
\section{Stage 1: Pot Bubble Point and Vapour Generation}
\label{sec:stage1}

\subsection{Bubble-point temperature}

The bubble-point temperature $T_1$ satisfies
\begin{equation}
  \mathcal{R}_{\mathrm{bp}}(T_1) \;=\; \sum_{i=1}^{n_c} K_i(T_1, P, \mathbf{X})\, X_i - 1 = 0.
  \label{eq:bpresidual}
\end{equation}
This is a monotone scalar equation in $T_1$ and is solved by \texttt{bisect}
on the bracket $[T_{\mathrm{lo}}, T_{\mathrm{hi}}]$ (default 375--700\,K),
with automatic bracket expansion in 10\,K steps if the initial bracket does
not straddle a sign change.

Once $T_1$ is found, the equilibrium vapour composition is
\begin{equation}
  y_{1,i} = K_i(T_1, P, \mathbf{X})\, X_i \bigg/ \sum_j K_j X_j.
\end{equation}

\subsection{Vapour generation rate}

The pot energy balance (quasi-steady, bubble-point liquid) gives the molar
vapour generation rate $D_1$\,[mol/s]:
\begin{equation}
  Q_1 + R_2\, C_{p,L}(T_1, \mathbf{X})\,(T_2 - T_1)
  \;=\; D_1\, \Delta H_{\mathrm{vap}}(T_1, \mathbf{X}),
  \label{eq:stage1energy}
\end{equation}
where
\begin{align}
  \Delta H_{\mathrm{vap}}(T, \mathbf{x}) &= \sum_i x_i\, M_i\, L_i(T) \quad [\J/\mol], \\
  C_{p,L}(T, \mathbf{x}) &= \sum_i x_i\, C_{p,L,i}(T) \quad [\J/\mol/\K],
\end{align}
with $L_i(T)$ the Watson-correlation latent heat [J/kg] and $C_{p,L,i}$ the
liquid heat capacity from group-contribution data.  Rearranging
\cref{eq:stage1energy}:
\begin{equation}
  D_1 = \frac{Q_1 + R_2\, C_{p,L}(T_2 - T_1)}{\Delta H_{\mathrm{vap}}(T_1, \mathbf{X})},
  \qquad D_1 \ge 0.
  \label{eq:D1}
\end{equation}

% ---------------------------------------------------------------------------
\section{Stage 2: Column Neck Partial Condensation}
\label{sec:stage2}

\subsection{Physical model}

The exposed column neck (length $L$, mean diameter $\bar{D}$, area $A = \pi \bar{D} L$)
loses heat to ambient air by natural convection.  The vapour stream $D_1$ at
temperature $T_1$ and composition $y_{1,i}$ enters the neck, partially
condenses, and splits into:
\begin{itemize}[nosep]
  \item reflux $R_2 = D_1(1-V_{\mathrm{fl}})$ [mol/s] returned to the pot,
  \item forward vapour $D_2 = D_1 V_{\mathrm{fl}}$ [mol/s] passing to the condenser.
\end{itemize}

\subsection{Natural-convection heat-transfer coefficient}

The Churchill--Chu correlation gives the Nusselt number for a vertical cylinder:
\begin{align}
  T_{\mathrm{film}} &= \tfrac{1}{2}(T_2 + T_{\mathrm{room}}),\\
  \mathrm{Ra} &= \mathrm{Gr}(T_2, T_{\mathrm{room}}, L) \cdot \mathrm{Pr}(T_{\mathrm{film}}),\\
  \mathrm{Nu} &= \begin{cases}
    0.68 + \dfrac{0.670\,\mathrm{Ra}^{1/4}}{\bigl[1+(0.492/\mathrm{Pr})^{9/16}\bigr]^{4/9}}
    & \mathrm{Ra} < 10^9,\\[8pt]
    \left(0.825 + \dfrac{0.387\,\mathrm{Ra}^{1/6}}{\bigl[1+(0.492/\mathrm{Pr})^{9/16}\bigr]^{8/27}}\right)^{\!2}
    & \mathrm{Ra} \ge 10^9,
  \end{cases}\\
  h &= h_{\mathrm{mult}} \cdot \mathrm{Nu}\, k_{\mathrm{air}}(T_{\mathrm{film}}) / L \quad [\W/\mathrm{m}^2/\K],
\end{align}
where $\mathrm{Gr}$, $\mathrm{Pr}$, and $k_{\mathrm{air}}$ are evaluated from
polynomial fits to air properties.

\subsection{Energy balance and neck temperature}

The neck-exit temperature $T_2$ satisfies the simultaneous energy balance:
\begin{equation}
  \underbrace{h A (T_2 - T_{\mathrm{room}})}_{\text{convective loss}}
  = \underbrace{D_2\, C_{p,v}(T_1 - T_2)
    + R_2\Bigl[\Delta H_{\mathrm{vap}}(T_2,\,\mathbf{x}_r) + C_{p,L}(T_1 - T_2)\Bigr]}_{\text{enthalpy released by cooling and condensation}},
  \label{eq:stage2energy}
\end{equation}
where $V_{\mathrm{fl}}(T_2)$, $\mathbf{x}_r$ (reflux composition), and $\mathbf{y}_2$
(forward vapour composition) are obtained from an isothermal Rachford--Rice
flash (\cref{sec:rr}) at $(T_2, P)$ with feed $z_i = y_{1,i}$.

\textbf{[Issue~6 fix]:} The latent heat $\Delta H_{\mathrm{vap}}$ on the
right-hand side of \cref{eq:stage2energy} is evaluated at the \emph{reflux
(liquid) composition} $\mathbf{x}_r$ from the flash, not at the feed
composition $\mathbf{y}_{1}$.  The condensed stream is enriched in heavier
components with higher $\Delta H_{\mathrm{vap}}$; using the feed composition
underestimates the latent heat released.

The scalar residual $Q_{\mathrm{loss}}(T_2) - Q_{\mathrm{supp}}(T_2)$ is
solved for $T_2 \in [T_{\mathrm{room}}+5, T_1]$ by \texttt{bisect}
(tolerance 0.1\,K).  Each bisect evaluation calls one Rachford--Rice flash,
hence this is the most expensive step per derivative evaluation.

% ---------------------------------------------------------------------------
\section{Stage 2b: Condenser}
\label{sec:condenser}

\subsection{Effective \texorpdfstring{$\varepsilon$}{e}-NTU with latent heat}

The D86 condenser is modelled as a single lumped heat-exchange element with
UA product $UA_{\mathrm{cond}}$\,[W/K] and an isothermal cooling bath at
$T_{\mathrm{bath}}$.  The vapour stream $D_2$ at $(T_2, \mathbf{y}_2)$ enters
and is partially condensed.

\textbf{[Issue~1 fix]:} The original model used the sensible vapour heat
capacity $C_{p,v}$ alone as the effective heat-capacity rate, ignoring
the latent heat of condensation.  The corrected effective rate is
\begin{equation}
  C_{\mathrm{eff}} = D_2 \left(C_{p,v}(T_2,\, \mathbf{y}_2)
    + \frac{\Delta H_{\mathrm{vap}}(T_2,\, \mathbf{y}_2)}{T_2 - T_{\mathrm{bath}}}\right)
  \quad [\W/\K].
  \label{eq:Ceff}
\end{equation}
The second term accounts for the latent heat released as vapour condenses
across the temperature drop $(T_2 - T_{\mathrm{bath}})$.

The NTU--effectiveness relations for an isothermal bath ($C_{\mathrm{max}} \to \infty$) give:
\begin{align}
  \mathrm{NTU} &= \frac{UA_{\mathrm{cond}}}{C_{\mathrm{eff}}}, \label{eq:NTU}\\
  \varepsilon   &= 1 - e^{-\mathrm{NTU}}, \label{eq:eps}\\
  T_{\mathrm{cond}} &= T_2 - \varepsilon\,(T_2 - T_{\mathrm{bath}}). \label{eq:Tcond}
\end{align}

One isothermal Rachford--Rice flash at $(T_{\mathrm{cond}}, P)$ with feed $\mathbf{y}_2$
determines the condenser vapour fraction $V_c$ and outlet compositions:
\begin{align}
  D_{c,\mathrm{vap}}  &= D_2\, V_c, \\
  D_{c,\mathrm{liq}}  &= D_2\,(1 - V_c), \\
  D_{c,\mathrm{tot}}  &= D_2.
\end{align}
A mole-weighted average condenser outlet composition is formed:
\begin{equation}
  \mathbf{c}_{\mathrm{cond}} = \frac{D_{c,\mathrm{vap}}\,\mathbf{y}_{c}
    + D_{c,\mathrm{liq}}\,\mathbf{x}_{c}}{D_{c,\mathrm{tot}}},
\end{equation}
normalised to unit sum.

% ---------------------------------------------------------------------------
\section{Stage 3: Thermocouple Thermal Lag}
\label{sec:stage3}

\subsection{Lumped-capacitance model with latent heat}

The D86 thermocouple bulb is modelled as a lumped thermal mass $C_{\mathrm{glass}}$\,[J/K]
driven by the vapour stream exiting the column neck.

\textbf{[Issue~5 fix]:} The original model drove the thermocouple only by
the sensible heat of the forward vapour $D_2$.  In the neck, the reflux
stream $R_2$ also condenses, releasing latent heat at the neck exit temperature.
This condensation contributes an additional heat-capacity rate to the bulb.
The effective heat-capacity rate is
\begin{equation}
  \dot{C}_{\mathrm{eff}} = \underbrace{D_2\, C_{p,v}(T_2,\,\mathbf{y}_2)}_{\text{sensible forward flow}}
  + \underbrace{\frac{R_2\, \Delta H_{\mathrm{vap}}(T_2,\,\mathbf{y}_2)}{T_1 - T_2}}_{\text{condensation latent heat}},
  \label{eq:Cflow}
\end{equation}
where the denominator $T_1 - T_2$ converts the latent heat flux to an
effective heat-capacity rate across the available temperature driving force.

\subsection{Analytical update}

The governing first-order ODE
\begin{equation}
  C_{\mathrm{glass}}\,\frac{\dd T_{\mathrm{D86}}}{\dd t}
  = \dot{C}_{\mathrm{eff}}\,(T_2 - T_{\mathrm{D86}})
\end{equation}
is solved analytically over one time step $\dt$:
\begin{equation}
  T_{\mathrm{D86}}^{n+1}
  = T_2 + \left(T_{\mathrm{D86}}^{n} - T_2\right)
    \exp\!\left(-\frac{\dot{C}_{\mathrm{eff}}\,\dt}{C_{\mathrm{glass}}}\right).
  \label{eq:Td86update}
\end{equation}
This solution is unconditionally stable and exact for constant $T_2$ and
$\dot{C}_{\mathrm{eff}}$ over the step.

% ---------------------------------------------------------------------------
\section{Pot Mole ODE}
\label{sec:ode}

The component moles in the pot evolve according to
\begin{equation}
  \frac{\dd N_i}{\dd t} = R_2\, x_{r,i} - D_1\, y_{1,i},
  \label{eq:ode}
\end{equation}
where $x_{r,i}$ is the reflux composition from Stage~2 and $y_{1,i}$ is the
vapour composition from Stage~1.  The condenser does not contribute reflux;
all condenser-outlet flow goes to the distillate receiver.

The right-hand side $\mathbf{f}(\mathbf{N}, T_2, R_2, Q_1) = (R_2 x_{r,i} - D_1 y_{1,i})$
is the derivative function \texttt{\_compute\_derivatives\_condenser\_v2},
which executes Stages~1--2b in sequence to evaluate all quantities.

% ---------------------------------------------------------------------------
\section{RK2 Time Integration (Heun's Method)}
\label{sec:rk2}

The mole ODE~\cref{eq:ode} is advanced by Heun's predictor--corrector:
\begin{align}
  \mathbf{k}_1 &= \mathbf{f}\!\left(\mathbf{N}^n,\; T_2^n,\; R_2^n,\; Q_1^n\right),
  \label{eq:k1}\\
  \mathbf{N}^{\mathrm{pred}} &= \max\!\left(\mathbf{N}^n + \dt\,\mathbf{k}_1,\, \mathbf{0}\right),
  \label{eq:pred}\\
  \mathbf{k}_2 &= \mathbf{f}\!\left(\mathbf{N}^{\mathrm{pred}},\; T_2^{(k_1)},\; R_2^{(k_1)},\; Q_1^n\right),
  \label{eq:k2}\\
  \mathbf{N}^{n+1} &= \max\!\left(\mathbf{N}^n + \frac{\dt}{2}(\mathbf{k}_1 + \mathbf{k}_2),\, \mathbf{0}\right),
  \label{eq:corrector}
\end{align}
where superscript $(k_1)$ denotes quantities evaluated during the $k_1$ pass
(used as the predictor state for $k_2$).  The positivity projection
$\max(\cdot, \mathbf{0})$ prevents small negative moles from floating-point
accumulation.

Averaged quantities for post-processing are formed as
\begin{equation}
  \bar{q} = \tfrac{1}{2}(q^{(k_1)} + q^{(k_2)}),
\end{equation}
and the Stage~3 thermocouple update (\cref{eq:Td86update}) uses $\bar{D}_2$,
$\bar{T}_2$, $\bar{R}_2$, $\bar{T}_1$, and $\bar{\mathbf{y}}_2$.

% ---------------------------------------------------------------------------
\section{Distillate Volume Accounting}
\label{sec:distvol}

\textbf{[Issue~7 fix]:} The original model counted only the condensed liquid
$D_{c,\mathrm{liq}}$ as distillate, losing the uncondensed vapour fraction
$D_{c,\mathrm{vap}}$ from the mass balance.  The corrected distillate molar
flow is the full condenser outlet:
\begin{equation}
  \Delta n_{\mathrm{dist}} = \bar{D}_{c,\mathrm{tot}}\,\dt
  = \bar{D}_2\,\dt \quad [\mol].
\end{equation}
The distillate volume increment at room temperature is
\begin{equation}
  \Delta V_{\mathrm{dist}} = \frac{\Delta n_{\mathrm{dist}}\,\bar{M}_{\mathrm{cond}}}
    {\rho_{\mathrm{cond}}(T_{\mathrm{room}})}\times 10^6 \quad [\mL],
  \label{eq:dV}
\end{equation}
where $\bar{M}_{\mathrm{cond}} = \sum_i c_{\mathrm{cond},i}\,M_i$ is the
condenser-outlet mean molar mass and $\rho_{\mathrm{cond}}$ is the liquid
density at room temperature.

The pot volume is updated from the accepted mole vector:
\begin{equation}
  V_{\mathrm{pot}} = \frac{W^{n+1}\,\bar{M}_{\mathrm{pot}}}
    {\rho_{\mathrm{pot}}(T_1^{(k_2)})} \times 10^6 \quad [\mL],
\end{equation}
with $\bar{M}_{\mathrm{pot}} = \sum_i X_i^{n+1} M_i$ evaluated at the
updated composition and $\rho_{\mathrm{pot}}$ at the Stage~1 bubble-point
temperature (liquid is at its bubble point).

% ---------------------------------------------------------------------------
\section{Option~E Controller}
\label{sec:controller}

\subsection{Inverse-model feed-forward}

Option~E inverts the Stage~1 energy balance~\cref{eq:stage1energy} to compute
the heater power $Q_1$ that \emph{exactly} achieves a target distillate molar
flow.

\textbf{[Issue~11 fix]:} The original controller used the pot composition
$\mathbf{X}$ to evaluate the liquid density for the volume-to-molar-flow
conversion.  Because the distillate is richer in lighter components, its
density differs from the pot liquid.  The corrected controller uses the
condenser-outlet composition $\mathbf{c}_{\mathrm{cond}}$ carried from the
previous step.

The target distillate molar flow is
\begin{equation}
  D_{\mathrm{dist,target}}
  = \frac{\rho_{\mathrm{cond}}(T_{\mathrm{room}})\,\dot{V}_{\mathrm{target}}}
    {60 \times 10^6 \times \bar{M}_{\mathrm{cond}}} \quad [\mol/\text{s}],
  \label{eq:Dtarget}
\end{equation}
where $\rho_{\mathrm{cond}}$ and $\bar{M}_{\mathrm{cond}}$ are evaluated at
the distillate composition $\mathbf{c}_{\mathrm{cond}}$.

The total vapour target from the pot must also supply the reflux that will
return from the neck:
\begin{equation}
  D_{1,\mathrm{target}} = D_{\mathrm{dist,target}} + R_2^n.
  \label{eq:D1target}
\end{equation}

Inverting \cref{eq:stage1energy} for $Q_1$ gives the feed-forward power:
\begin{equation}
  Q_1^{\mathrm{ff}} = D_{1,\mathrm{target}}\,\Delta H_{\mathrm{vap}}(T_1^{(k_2)}, \mathbf{X})
    + R_2^n\,C_{p,L}(T_1^{(k_2)}, \mathbf{X})\,(T_1^{(k_2)} - T_2^n),
  \quad Q_1^{\mathrm{ff}} \ge 0.
  \label{eq:Q1ff}
\end{equation}

\subsection{Rate-tracking proportional trim}

An EMA-smoothed instantaneous distillation rate is maintained:
\begin{equation}
  \dot{V}_{\mathrm{EMA}}^{n+1}
  = \alpha\,\frac{\Delta V_{\mathrm{dist}}}{\dt}\times 60
    + (1 - \alpha)\,\dot{V}_{\mathrm{EMA}}^{n}
  \quad [\mL/\min],
\end{equation}
with smoothing coefficient $\alpha = 0.15$.  A proportional trim corrects
any residual rate error:
\begin{equation}
  Q_1^{\mathrm{trim}} = K_{p,\mathrm{trim}}\,(\dot{V}_{\mathrm{target}} - \dot{V}_{\mathrm{EMA}}^{n+1})
  \quad [\W].
\end{equation}

The commanded heater power is
\begin{equation}
  Q_1^{n+1} = \mathrm{clip}\!\left(Q_1^{\mathrm{ff}} + Q_1^{\mathrm{trim}},\;
    Q_{\min},\; Q_{\max}\right).
  \label{eq:Q1}
\end{equation}
Because the feed-forward handles the dominant physics, $K_{p,\mathrm{trim}}$
can be kept small (default 10\,W\,min/mL) and the trim rarely saturates.

% ---------------------------------------------------------------------------
\section{Stall Detection}
\label{sec:stall}

Two stall conditions abort the simulation early:

\paragraph{Pre-drop stall.}
If no distillate has been collected ($V_{\mathrm{dist}} \le \epsilon_V$)
\emph{and} the forward vapour flow is negligible ($\bar{D}_2 \le \epsilon_D$)
for longer than $\tau_{\mathrm{startup}}$ seconds of simulated time, the
simulation terminates.  This catches parameter sets where the heater cannot
maintain boiling.

\paragraph{Post-drop stall.}
After the first drop is collected, the stall timer resets whenever the
distillate volume advances by more than $\epsilon_V$.  If less than $\epsilon_V$
accumulates in $\tau_{\mathrm{stall}}$ seconds, the simulation terminates.
The threshold $\epsilon_V = 0.5$\,mL prevents near-zero-rate runs (e.g.,
caused by excessive condenser duty) from running indefinitely.

Default values: $\epsilon_V = 0.5$\,mL, $\epsilon_D = 10^{-9}$\,mol/s,
$\tau_{\mathrm{stall}} = 300$\,s, $\tau_{\mathrm{startup}} = 900$\,s.

% ---------------------------------------------------------------------------
\section{Complete Algorithm}
\label{sec:algorithm}

\begin{algorithm}[H]
\caption{D86 simulation — \texttt{run\_d86\_simulation\_rk2\_condenser\_optE\_v2}}
\begin{algorithmic}[1]
\State \textbf{Initialise} $\mathbf{N}^0$, $T_2^0 = T_{\mathrm{room}}$, $R_2^0 = 0$,
       $T_{\mathrm{D86}}^0 = T_{\mathrm{room}}$, $Q_1^0$, $\mathbf{c}_{\mathrm{cond}}^0 = \mathbf{X}^0$
\While{$V_{\mathrm{pot}} > V_{\min}$ \textbf{and} $W > 0$}
  \State \textbf{--- $k_1$ pass ---}
  \State $T_1^{(1)}, \mathbf{y}_1^{(1)} \leftarrow$ \Call{BubblePoint}{$\mathbf{X}^n, P$} \hfill\Comment{\cref{eq:bpresidual}}
  \State $D_1^{(1)} \leftarrow$ \Call{EnergyBalance}{$T_1^{(1)}, Q_1^n, R_2^n, T_2^n$} \hfill\Comment{\cref{eq:D1}}
  \State $T_2^{(1)}, R_2^{(1)}, \mathbf{x}_r^{(1)}, D_2^{(1)}, \mathbf{y}_2^{(1)} \leftarrow$ \Call{NeckFlash}{$D_1^{(1)}, T_1^{(1)}, \mathbf{y}_1^{(1)}$} \hfill\Comment{\cref{eq:stage2energy}}
  \State $T_c^{(1)}, D_{cv}^{(1)}, D_{cl}^{(1)}, \mathbf{c}^{(1)} \leftarrow$ \Call{Condenser}{$D_2^{(1)}, T_2^{(1)}, \mathbf{y}_2^{(1)}$} \hfill\Comment{\cref{eq:Ceff}--\cref{eq:Tcond}}
  \State $\mathbf{k}_1 \leftarrow R_2^{(1)}\mathbf{x}_r^{(1)} - D_1^{(1)}\mathbf{y}_1^{(1)}$ \hfill\Comment{\cref{eq:ode}}
  \State $\mathbf{N}^{\mathrm{pred}} \leftarrow \max(\mathbf{N}^n + \dt\,\mathbf{k}_1,\, \mathbf{0})$ \hfill\Comment{\cref{eq:pred}}
  \State \textbf{--- $k_2$ pass (uses $\mathbf{N}^{\mathrm{pred}}$, $T_2^{(1)}$, $R_2^{(1)}$) ---}
  \State $T_1^{(2)}, \mathbf{y}_1^{(2)}, D_1^{(2)}, T_2^{(2)}, R_2^{(2)}, D_2^{(2)}, \mathbf{y}_2^{(2)}, \mathbf{c}^{(2)} \leftarrow \ldots$ (same sequence)
  \State $\mathbf{k}_2 \leftarrow R_2^{(2)}\mathbf{x}_r^{(2)} - D_1^{(2)}\mathbf{y}_1^{(2)}$
  \State $\mathbf{N}^{n+1} \leftarrow \max\!\bigl(\mathbf{N}^n + \tfrac{\dt}{2}(\mathbf{k}_1+\mathbf{k}_2),\, \mathbf{0}\bigr)$ \hfill\Comment{\cref{eq:corrector}}
  \State \textbf{--- Averages ---}
  \State $\bar{D}_2, \bar{T}_2, \bar{R}_2, \bar{T}_1, \bar{\mathbf{y}}_2, \bar{\mathbf{c}}_{\mathrm{cond}} \leftarrow \tfrac{1}{2}(\cdot^{(1)} + \cdot^{(2)})$
  \State \textbf{--- Stage 3: thermocouple update ---}
  \State $\dot{C}_{\mathrm{eff}} \leftarrow \bar{D}_2 C_{p,v}(\bar{T}_2, \bar{\mathbf{y}}_2) + \bar{R}_2 \Delta H_{\mathrm{vap}}(\bar{T}_2, \bar{\mathbf{y}}_2)/(\bar{T}_1 - \bar{T}_2)$ \hfill\Comment{\cref{eq:Cflow}}
  \State $T_{\mathrm{D86}}^{n+1} \leftarrow \bar{T}_2 + (T_{\mathrm{D86}}^n - \bar{T}_2)\exp(-\dot{C}_{\mathrm{eff}}\,\dt / C_{\mathrm{glass}})$ \hfill\Comment{\cref{eq:Td86update}}
  \State \textbf{--- Accept state ---}
  \State $R_2^{n+1} \leftarrow \bar{R}_2$; \quad $T_2^{n+1} \leftarrow \bar{T}_2$; \quad $\mathbf{c}_{\mathrm{cond}}^{n+1} \leftarrow \bar{\mathbf{c}}_{\mathrm{cond}}$
  \State \textbf{--- Distillate and pot volumes ---}
  \State $\Delta V_{\mathrm{dist}} \leftarrow \bar{D}_2\,\dt\,\bar{M}_{\mathrm{cond}} / \rho_{\mathrm{cond}}(T_{\mathrm{room}}) \times 10^6$ \hfill\Comment{\cref{eq:dV}}
  \State $V_{\mathrm{dist}} \mathrel{+}= \Delta V_{\mathrm{dist}}$; \quad update $V_{\mathrm{pot}}$
  \State \textbf{--- Option E controller ---}
  \State Compute $Q_1^{\mathrm{ff}}$ from \cref{eq:Q1ff} using $\mathbf{c}_{\mathrm{cond}}^{n+1}$ for density \hfill\Comment{Issue~11}
  \State Update $\dot{V}_{\mathrm{EMA}}$; compute $Q_1^{n+1}$ from \cref{eq:Q1}
  \State \textbf{--- Stall check ---}; \textbf{break} if stalled \hfill\Comment{\cref{sec:stall}}
  \State $t \mathrel{+}= \dt$
\EndWhile
\end{algorithmic}
\end{algorithm}

% ---------------------------------------------------------------------------
\section{Summary of Physics Corrections}
\label{sec:fixes}

\begin{center}
\renewcommand{\arraystretch}{1.4}
\begin{tabular}{cllp{5cm}}
\toprule
Issue & Location & Symptom & Correction \\
\midrule
1 & Stage~2b condenser & Underestimated heat removal; too little condensation & Add $L_v/(T_{\mathrm{in}}{-}T_{\mathrm{bath}})$ to effective $C_{\mathrm{eff}}$ (\cref{eq:Ceff}) \\
5 & Stage~3 thermocouple & Thermocouple too slow; misses condensation latent heat & Add $R_2 L_v/(T_1{-}T_2)$ to $\dot{C}_{\mathrm{eff}}$ (\cref{eq:Cflow}) \\
6 & Stage~2 neck flash & Lv underestimated for heavy-enriched condensate & Evaluate $\Delta H_{\mathrm{vap}}$ at reflux composition $\mathbf{x}_r$ \\
7 & Distillate accounting & Uncondensed vapour fraction lost from mass balance & Count $D_{c,\mathrm{tot}} = D_2$ not $D_{c,\mathrm{liq}}$ (\cref{eq:dV}) \\
11 & Option E controller & Density reference at wrong (pot) composition & Use $\mathbf{c}_{\mathrm{cond}}$ for $\rho$ and $\bar{M}$ in \cref{eq:Dtarget} \\
\bottomrule
\end{tabular}
\end{center}

% ---------------------------------------------------------------------------
\section{Simulation Parameters Reference}
\label{sec:params}

\begin{center}
\begin{tabular}{llll}
\toprule
Parameter & Symbol & Default & Description \\
\midrule
\texttt{dt} & $\dt$ & 2.0\,s & Integration time step \\
\texttt{UA\_cond} & $UA_{\mathrm{cond}}$ & calibrated & Condenser UA product (W/K) \\
\texttt{T\_bath} & $T_{\mathrm{bath}}$ & calibrated & Cooling-bath temperature (K) \\
\texttt{h\_multiplier} & $h_{\mathrm{mult}}$ & calibrated & Neck convection scaling \\
\texttt{C\_glass} & $C_{\mathrm{glass}}$ & calibrated & Thermocouple heat capacity (J/K) \\
\texttt{T\_bubble\_lo/hi} & $T_{\mathrm{lo}}, T_{\mathrm{hi}}$ & 375, 700\,K & Bubble-point search bounds \\
\texttt{target\_rate\_ml\_min} & $\dot{V}_{\mathrm{target}}$ & 4.5\,mL/min & Target distillation rate \\
\texttt{controller\_Kp\_trim} & $K_{p,\mathrm{trim}}$ & 10\,W\,min/mL & Proportional trim gain \\
\texttt{Q1\_min/max} & $Q_{\min}, Q_{\max}$ & 0, 1500\,W & Heater power limits \\
\texttt{rate\_ema\_alpha} & $\alpha$ & 0.15 & Rate EMA smoothing \\
\texttt{stall\_vol\_tol\_mL} & $\epsilon_V$ & 0.5\,mL & Stall volume threshold \\
\texttt{stall\_window\_s} & $\tau_{\mathrm{stall}}$ & 300\,s & Post-drop stall window \\
\texttt{startup\_stall\_window\_s} & $\tau_{\mathrm{startup}}$ & 900\,s & Pre-drop stall window \\
\texttt{use\_srk} & --- & False & Use SRK EoS (vs.\ UNIFAC) \\
\bottomrule
\end{tabular}
\end{center}

\end{document}
